\FSabschnitt{Teil A \quad Alle Erklärfragen — an einem Platz}{Klausurteil A1 / A2}



% =====================================================================
\begin{RezeptA}{2.3\ \ Halbleiter, Dotierung, Leitung \hfill
  \Quelle{Ü1 A2 · Ü4 A3 · PK13 A1 · Kl09 · Kl10 · Kl11 (6×)}}
\Erkennung{„Was versteht man unter …“ $\cdot$ „Erklären Sie …“
$\cdot$ Ankreuzfragen zu Halbleitereigenschaften.}
\FSline
\FA{Was ist ein Halbleiter?}{Ein Festkörper, dessen Leitfähigkeit
\erg{zwischen der eines Metalls und der eines Isolators} liegt. Seine
Bandlücke ist klein genug, dass thermische Energie Elektron-Loch-Paare
erzeugt.}

\FA{Wie ändert sich $\sigma$ mit der Temperatur?}{Beim Halbleiter
\erg{nimmt $\sigma$ mit fallender Temperatur ab}, weil weniger Paare
thermisch erzeugt werden. Beim Metall ist es umgekehrt: dort steigt
$\sigma$ bei Abkühlung, weil die Gitterschwingungen abnehmen.}

\FA{Was ist Eigenleitung?}{Die Leitung im \emph{un}dotierten Halbleiter,
die allein von thermisch erzeugten Elektron-Loch-Paaren getragen wird.
Dabei gilt $n=p=n_i$.}

\FA{Was ist Paarbildung?}{Thermische oder optische Energie hebt ein
Elektron aus dem Valenzband ins Leitungsband. Dabei entsteht
\erg{gleichzeitig} ein freies Elektron und ein Loch — immer paarweise.}

\FA{Was ist Rekombination?}{Ein freies Elektron fällt zurück in ein Loch,
das Paar verschwindet. Die frei werdende Energie wird als Wärme oder als
Photon abgegeben.}

\FA{Was heißt thermodynamisches Gleichgewicht?}{Paarbildung und
Rekombination laufen gleich schnell ab, die Trägerdichten sind zeitlich
konstant. Dann gilt \erg{$n\cdot p=n_i^{2}$}.}

\FA{Was bewirkt Dotieren?}{Das gezielte Einbringen von Fremdatomen,
um \emph{eine} Ladungsträgersorte um Größenordnungen zu erhöhen und
damit die Leitfähigkeit einzustellen.}

\FA{Wodurch entsteht n-Leitung?}{Durch fünfwertige Dotanden
(\erg{Phosphor, Arsen}). Vier Elektronen gehen Bindungen ein, das fünfte
ist nur schwach gebunden und wird schon bei Raumtemperatur frei.
Es gilt $N_D>N_A$, Elektronen sind Majoritäten.}

\FA{Wodurch entsteht p-Leitung?}{Durch dreiwertige Dotanden
(\erg{Bor, Aluminium, Indium}). Ihnen fehlt ein Bindungselektron, die
Lücke wirkt als Loch. Es gilt $N_A>N_D$, Löcher sind Majoritäten.
\warn{Silizium selbst ist vierwertig und Grundmaterial, kein Dotand.}}

\FA{Majoritäts- und Minoritätsträger?}{Majoritäten sind die häufigere
Sorte (im n-HL die Elektronen, $n_0\approx N_D$). Minoritäten sind die
seltenere Sorte, ihre Dichte folgt aus \erg{$p_0=n_i^{2}/n_0$}.}

\FA{Warum ist ein n-Halbleiter nach außen neutral?}{Weil zu jedem
abgegebenen freien Elektron ein \erg{ortsfester, positiv geladener
Donatorrumpf} im Gitter zurückbleibt. Beide Ladungen sind gleich groß
und heben sich in der Summe auf.}

\FA{Was ist Störstellenerschöpfung?}{Der Temperaturbereich, in dem
\emph{alle} Dotanden ionisiert sind, die Eigenleitung aber noch
vernachlässigbar ist. Dort ist $n_0\approx N_D$ konstant — das ist der
\erg{Arbeitsbereich} des Bauelements.}

\FA{Drift oder Diffusion?}{\textbf{Drift} ist Bewegung \emph{im
elektrischen Feld}: $j=\sigma E$. \textbf{Diffusion} ist Bewegung
\emph{entlang eines Konzentrationsgefälles}: $j=-e\,D\,\dd n/\dd x$.}

\FA{Was sagt die Kontinuitätsgleichung?}{Sie bilanziert die
Trägerdichte: zeitliche Änderung $=$ Zufluss $-$ Abfluss $+$ Generation
$-$ Rekombination.}
\end{RezeptA}

% =====================================================================
\begin{RezeptA}{2.4\ \ Bauelemente in einem Satz \hfill
  \Quelle{PK-I A1 · PK13 A1 · BA A1–A5 · Kl11 · Kl12 (5×)}}
\Erkennung{„Was ist eine …?“ $\cdot$ „Wofür steht die Abkürzung …?“
$\cdot$ „Nennen Sie das Funktionsprinzip.“}
\FSline
\FA{Was ist eine Schottky-Diode?}{Ein \erg{Metall-Halbleiter-Übergang}
statt eines pn-Übergangs. Sie hat eine kleinere Flussspannung und
schaltet schneller, weil keine Diffusionsladung gespeichert wird.}

\FA{Wofür steht CMOS?}{\erg{Complementary Metal Oxide Semiconductor} —
je ein selbstsperrender n- und p-Kanal-MOSFET, die komplementär
schalten. Im statischen Zustand fließt fast kein Querstrom.}

\FA{Wie funktioniert eine Solarzelle?}{Über den \erg{inneren
Fotoeffekt am pn-Übergang}: Photonen mit $W>W_g$ erzeugen Paare in der
Raumladungszone, das eingebaute Feld trennt sie und treibt einen Strom.}

\FA{Was heißt Gleichrichten?}{Strom nur in \emph{eine} Richtung
durchlassen und so eine Wechselgröße in eine Gleichgröße umwandeln —
die Aufgabe der Diode.}

\FA{Warum steuert ein FET leistungslos, ein BJT nicht?}{Der FET steuert
über das \erg{elektrische Feld} am isolierten Gate; im stationären
Betrieb fließt kein Gatestrom, also wird keine Steuerleistung gebraucht.
Der Bipolartransistor wird über den \erg{Basisstrom} gesteuert, und
dieser Strom kostet Leistung.}

\FA{Sind die Energien im unendlich hohen Potentialtopf gequantelt oder
kontinuierlich?}{\erg{Gequantelt} — es sind nur diskrete Eigenwerte
$E_n\propto n^{2}/L^{2}$ erlaubt. Je größer $L$, desto kleiner der
Abstand $\Delta E$, bis er im Festkörper quasi-kontinuierlich wird.}
\end{RezeptA}

% =====================================================================
\begin{RezeptA}{2.5\ \ Prinzip der Gegenkopplung \hfill
  \Quelle{alle 6 Altklausuren · Ü8 A4 · BA B8 \ \erg{(6/6 — sicher dran)}}}
\Erkennung{„Erläutern Sie das Prinzip der Gegenkopplung.“ $\cdot$
„Welche Vor- und Nachteile hat sie?“ $\cdot$ LED mit
Konstantstromquelle im Schaltbild.}
\FSline
\FA{a) Welche Vorteile hat die Gegenkopplung?}{Der Arbeitspunkt wird
gegen Temperaturgang und Exemplarstreuung \erg{stabilisiert}. Zusätzlich
sinken die Verzerrungen, die Bandbreite steigt, und Ein- und
Ausgangswiderstand werden definiert einstellbar.}

\FA{b) Welchen Nachteil hat sie?}{\erg{Die Verstärkung sinkt.} Man
tauscht Verstärkung gegen Stabilität.}

\FA{c) Wie wirkt sie in der LED-Schaltung?}{Als Regelkreis, Schritt für
Schritt: $I_{LED}$ steigt $\to$ Spannung über $R_2$ steigt $\to$ $U_{BE}$
von T2 steigt $\to$ T2 leitet stärker $\to$ T2 zieht Basisstrom von T1 ab
$\to$ T1 sperrt mehr $\to$ \erg{$I_{LED}$ sinkt wieder}. Die Störung
wird also ihrer eigenen Ursache entgegengesetzt.}
\end{RezeptA}


